%%%%%%%%%%%%%%%%%%%%%%%%%%%%%%%%%%%%%%%%%%%%%%%%%%%%%%%%%%%%%%%%%%%%%%%
%                                                                     %
%                        Charte du club Bar                           %
%                                                                Beta %
%%%%%%%%%%%%%%%%%%%%%%%%%%%%%%%%%%%%%%%%%%%%%%%%%%%%%%%%%%%%%%%%%%%%%%%
%
%  Auteur :
%    Julien Déoux <julien.deoux@telecomnancy.net>
%    
%  Licence :
%    CC-BY-NC 4.0 (http://creativecommons.org/licenses/by-nc/4.0)
%
%%%



%%%%%%%%%%%%%%%%%%%
%  Configuration  %
%%%%%%%%%%%%%%%%%%%

\documentclass{article} % Jusqu'à preuve du contraire, les documents édités par le
                                                     % Ceten ne font pas 300 pages
\usepackage[a4paper,includeheadfoot,margin=2.54cm]{geometry}
\usepackage[francais]{babel}

\usepackage{hyperref}
\usepackage{graphicx} 
\usepackage{titlesec} 
\usepackage{ifxetex} 
\usepackage[usenames]{xcolor} 
\definecolor{newCeten}{RGB}{130,11,95}
\usepackage{enumitem,amssymb}
\newlist{todolist}{itemize}{2}
\setlist[todolist]{label=$\square$}
\usepackage{multicol}
\usepackage{setspace}
\usepackage{fontspec} 
\usepackage[all]{nowidow}

\setmainfont{Roboto Slab}
\setsansfont{Roboto}
\setmonofont{Roboto Mono}
\newfontfamily\condensed{Roboto Condensed}
\newfontfamily\condensedlight{Roboto Condensed Light}
\newfontfamily\light{Roboto Slab Light}
\titleformat*{\section}{\large\condensedlight\color{newCeten}}

\title{Charte du club Bar}
\author{Julien Déoux}
\date\today

%%%%%%%%%%%%%%
%  Document  %
%%%%%%%%%%%%%%

\begin{document}
	
	%---------------%
	% Page de garde %
	%---------------%

	\begin{titlepage}
		\begin{center}
			\includegraphics[width=\textwidth]{images/ceten}\par
			\vspace{2cm}
			{\Huge \light{} Charte du club bar}\par
			\vfill
			\begin{spacing}{1.5}
				année: \underline{\hspace{2cm}}\\
				président: \underline{\hspace{8cm}}\\
				trésorier: \underline{\hspace{8cm}}\\
				secrétaire: \underline{\hspace{8cm}}\\
				vice-président: \underline{\hspace{8cm}}\\
				responsable logistique: \underline{\hspace{8cm}}\\
				responsable informatique: \underline{\hspace{8cm}}\\
				référent de TELECOM Nancy: \underline{\hspace{8cm}}\\
			\end{spacing}
			\vspace{\baselineskip}
			objet du club: Tenir le bar de l'école et vendre des produits
			alimentaires aux élèves
			\vfill
			{\footnotesize \light{} Ce texte régit le fonctionnement, les
			obligations et les droits attribués au club bar, club du Ceten comme
			défini dans les statuts. Un exemplaire est remis au président du
			club, l’autre est archivé au sein du BDE\@. \\
			Charte modifiée et votée le 12/03/2013 par le BDE du Ceten, avec
			l'accord du président du club bar}
		\end{center}
	\end{titlepage}

	%----------%
	% Articles %
	%----------%

	\pagenumbering{arabic}

	\begin{multicols}{2}

		\section{Constitution}
\label{sec:constitution}
		
		{\small

			Le club bar est une entité interne au Ceten, composé de membres et
			dédié à l'objet défini au début de cette charte. Le fonctionnement
			et les engagements pris par ce club sont définis dans la présente
			charte.

			L'existence de ce club est subordonné à la signature de la présente
			charte par le responsable des clubs du BDE, le référent de TELECOM
			Nancy et par les président, trésorier, secrétaire, vice-président,
			responsable logistique et responsable informatique du club.
			
			Faisant partie intégrante de l’association, toutes les règles et
			décisions s’imposant au Ceten s’appliquent également au club.

		}

		\section{Composition}
\label{sec:composition}

		{\small
		
			Le club bar est composé de membres qui doivent également être
			membres du Ceten. Ce club se compose au minimum de six personnes
			distinctes qui sont:
			\begin{itemize}
				\item le président, dont le rôle est défini en
					Section~\ref{sec:presidence}
				\item le trésorier, dont le rôle est défini en
					Section~\ref{sec:tresorerie}
				\item le secrétaire, dont le rôle est défini en
					Section~\ref{sec:secretariat}
				\item le vice-président, qui soutient le président dans ses
					fonctions
				\item le responsable logistique, qui surveille l'état des stocks
					et le bon déroulement des approvisionnements du club
				\item le responsable informatique, qui s'assure du bon
					fonctionnement des matériels et logiciels informatiques
					nécessaires au fonctionnement du club
			\end{itemize}

			Ces six personnes déterminent le «bureau restreint» du club, qui a
			pour fonction de gérer l'ensemble du club. Le club n'a lieu
			d'exister sans la présence de ces six personnes.

			La structure interne du club et laissée libre au bureau restreint.
			En particulier, il peut désigner un «bureau complémentaire» dont le
			rôle est d'aider le bureau restreint à assurer ses fonctions, et
			dont la constitution doit être immédiatement transmise au
			responsable des clubs du BDE\@.

			L'ensemble de ces deux bureaux constitue le «bureau étendu» du club.
			Dans la suite de cette charte, l'appellation «bureau» désigne le
			bureau étendu du club.

		}

		\section{Durée}
\label{sec:duree}

		{\small
		
			Un mandat du club bar commence au 1\up{er} janvier et se termine au
			31 décembre de l'année civile spécifiée au début de la charte, sauf
			en cas de réélection d'un ou plusieurs membres du bureau restreint
			du club.

			Un membre du club l'est de la date de son inscription jusqu'au
			1\up{er} septembre suivant.

		}

		\section{Présidence}
\label{sec:presidence}

		\section{Trésorerie}
\label{sec:tresorerie}

		\section{Secrétariat}
\label{sec:secretariat}

	\end{multicols}

\end{document}
